% Part: first-order-logic
% Chapter: natural-deduction
% Section: quantifier-rules

\documentclass[../../../include/open-logic-section]{subfiles}

\begin{document}

\olfileid{fol}{ntd}{qrl}

\olsection{পরিমাণসূচক-বিধি}

\subsection{$\lforall$-এর বিধি}

\begin{defish}
\AxiomC{$!A(a)$}
\RightLabel{\Intro{\lforall}}
\UnaryInfC{$\lforall[x][\Atom{!A}{x}]$}
\DisplayProof
\hfill
\AxiomC{$\lforall[x][\Atom{!A}{x}]$}
\RightLabel{\Elim{\lforall}}
\UnaryInfC{$!A(t)$}
\DisplayProof
\end{defish}

$\lforall$-এর বিধিগুলিতে $t$ একটি বদ্ধ পদ (অর্থাৎ কোনো চলরাশি নেই
এমন পদ), আর $a$ এমন !!a{constant} যা সিদ্ধান্ত
$\lforall[x][!A(x)]$-এ অথবা পূর্বধারণা~$!A(a)$-তে শেষ হওয়া
!!{derivation}-এর কোনো !!{undischarged} অনুমিতিতে আসে না।
\Intro{\lforall} অনুমানের~$a$-কে তার \emph{আইগেনচল} বলি।
\footnote{উপরের বিধিতে $a$ একটি ধ্রুবক হলেও আমরা ``আইগেনচল'' শব্দটি
ব্যবহার করি। এর কারণ ঐতিহাসিক।}

\subsection{$\lexists$-এর বিধি}

\begin{defish}
\AxiomC{$\Atom{!A}{t}$}
\RightLabel{\Intro{\lexists}}
\UnaryInfC{$\lexists[x][\Atom{!A}{x}]$}
\DisplayProof
\hfill
\AxiomC{$\lexists[x][\Atom{!A}{x}]$}
\AxiomC{[$\Atom{!A}{a}$]$^n$}
\DeduceC{$!C$}
\DischargeRule{\Elim{\lexists}}{n}
\BinaryInfC{$!C$}
\DisplayProof
\end{defish}

এখানেও $t$ একটি বদ্ধ পদ, আর $a$ এমন !!a{constant} যা পূর্বধারণা
$\lexists[x][!A(x)]$-এ, সিদ্ধান্ত~$!C$-তে অথবা দুই পূর্বধারণায় শেষ
হওয়া !!{derivation}s-এর কোনো !!{undischarged} অনুমিতিতে আসে না
($!A(a)$ অনুমিতিগুলি ছাড়া)। \Elim{\lexists} অনুমানের~$a$-কে তার
\emph{আইগেনচল} বলি।

\Intro{\lforall} অনুমানে আইগেনচলটি সিদ্ধান্তে বা পূর্বধারণার দিকে
নিয়ে যাওয়া সংশ্লিষ্ট নিষ্পাদনের কোনো !!{undischarged} অনুমিতিতে আসে না।
\Elim{\lexists} অনুমানে সেটি প্রধান অস্তিত্বসূচক পূর্বধারণা, সিদ্ধান্ত
বা পূর্বধারণাগুলির দিকে নিয়ে যাওয়া !!{derivation}s-এর কোনো অনবমুক্ত
অনুমিতিতে আসে না—বিশেষ অবমুক্তযোগ্য অনুমিতিগুলি এর ব্যতিক্রম। এই
বিধি-নির্দিষ্ট সীমাগুলিকে \emph{আইগেনচল-শর্ত} বলা হয়।
% BN-SRC-034 TARGET-CORRECTION-BEGIN
\par\smallskip\noindent\textnormal{\small\textbf{উৎস-সংশোধন (BN-SRC-034):} হিমায়িত ইংরেজি সারাংশে আইগেনচলটি “premises”-এও না-আসার কথা বলা হয়েছে, অথচ প্রদর্শিত $\Intro{\lforall}$ ও $\Elim{\lexists}$—উভয় বিধির অবমুক্তযোগ্য পূর্বধারণায় আইগেনচলটি অবশ্যই আসে। বাংলা সারাংশে আগের দুটি সুনির্দিষ্ট শর্তই পুনর্ব্যক্ত করা হয়েছে: সার্বিক প্রবর্তনে সিদ্ধান্ত ও অনবমুক্ত অনুমিতি, আর অস্তিত্বসূচক অপসারণে প্রধান পূর্বধারণা, সিদ্ধান্ত ও বিশেষ অনুমিতি-ব্যতীত অনবমুক্ত অনুমিতিতে আইগেনচলটি অনুপস্থিত থাকবে।} % BN-SRC-034 TARGET-CORRECTION-END

\begin{explain}
মনে করো, যখন !!{variable}~$x$ মুক্ত এমন কোনো !!a{formula}~$!A$ থাকে, তখন আমরা
$!A(x)$ লিখে তা দেখাই। একই প্রসঙ্গে $!A(t)$ হলো~$\Subst{!A}{t}{x}$-এর
সংক্ষিপ্ত রূপ। তাই $\Intro\lexists$ বিধিটি এভাবেও লেখা যেত:
\begin{prooftree}
\AxiomC{$\Subst{!A}{t}{x}$}
\RightLabel{\Intro{\lexists}}
\UnaryInfC{$\lexists[x][!A]$}
\end{prooftree}
লক্ষ করো, $t$ আগে থেকেই~$!A$-তে থাকতে পারে; যেমন, $!A$ হতে পারে
$\Atom{\Obj P}{t,x}$। তাই~$\Atom{\Obj P}{t,t}$ থেকে
$\lexists[x][\Atom{\Obj P}{t,x}]$ অনুমান করা $\Intro\lexists$-এর
সঠিক প্রয়োগ—পূর্বধারণায়~$t$-র এক বা একাধিক আবির্ভাব, সবগুলি
অবশ্যই নয়, বদ্ধ !!{variable}~$x$ দিয়ে ``প্রতিস্থাপন'' করা যায়। কিন্তু
$\Intro\lforall$ ও~$\Elim\lexists$-এর আইগেনচল-শর্তে !!{constant}~$a$
যেন~$!A$-তে না আসে। তাই $\Atom{\Obj P}{a,a}$ থেকে
$\lforall[x][\Atom{\Obj P}{a,x}]$ অনুমান করতে $\Intro\lforall$-এর
সঠিক প্রয়োগ করা যায় না।
\end{explain}

\begin{explain}
\Intro{\lexists} ও \Elim{\lforall}-এ বদ্ধ পদ~$t$-র ওপর আর কোনো
বিধিনিষেধ নেই, তাই কোনো অতিরিক্ত শর্ত নিয়ে ভাবতে হয় না। অন্যদিকে,
\Elim{\lexists} ও \Intro{\lforall} বিধিতে আইগেনচল-শর্ত অনুযায়ী
!!{constant}~$a$ সিদ্ধান্তে বা কোনো !!{undischarged} অনুমিতিতে কোথাও
আসতে পারবে না। পদ্ধতিটি যে বিশুদ্ধ—অর্থাৎ যে !!{undischarged}
অনুমিতিগুলি থেকে বাক্যগুলি অনুসৃত হয়, শুধু সেগুলি থেকেই !!{sentence}s
!!{derive}s করে—তা নিশ্চিত করতে এই শর্ত দরকার। এই শর্ত না থাকলে
নিচেরটি অনুমোদিত হতো:
\begin{prooftree}
  \AxiomC{$\lexists[x][!A(x)]$}
  \AxiomC{$\Discharge{!A(a)}{1}$}
  \RightLabel{*\Intro{\lforall}}
  \UnaryInfC{$\lforall[x][!A(x)]$}
  \RightLabel{\Elim{\lexists}}
  \BinaryInfC{$\lforall[x][!A(x)]$}
\end{prooftree}
কিন্তু $\lexists[x][!A(x)] \Entails/ \lforall[x][!A(x)]$।

পরিমাণসূচকের অপসারণ-বিধিগুলি !!{variable}s-এর জায়গায় শুধু বদ্ধ পদ
বসাতে দেয়; তাই !!{sentence}s-এর কোনো সেট থেকে নিষ্পন্ন হতে পারে এমন
প্রত্যেক !!{formula} নিজেও !!a{sentence}।
\end{explain}


\end{document}
